\documentclass[12pt,a4paper]{article}
\usepackage[utf8]{inputenc}
\usepackage[T2A]{fontenc}
\usepackage[russian]{babel}
\usepackage{amsmath,amssymb,amsthm}
\usepackage{booktabs,array}
\usepackage{geometry}
\usepackage{microtype}
\geometry{margin=2.3cm}

\newtheorem{definition}{Определение}
\newtheorem{proposition}{Предложение}
\newtheorem{protocol}{Протокол}
\newtheorem{nogocriterion}{Критерий провала}
\DeclareMathOperator{\Tr}{Tr}

\title{Оператор спектральных адресов $\pi$}
\author{}
\date{4 августа 2026 г.}

\begin{document}
\maketitle

\section{От набора формул к одной алгебре}

Рассмотрим одиннадцать коротких формул RPFT-атласа: $\alpha_s$,
$\sin^2\theta_W$, отношения масс лёгких кварков, ядро отношения
$m_\tau/m_\mu$, $V_{cb}$ и три космологических доли. Прямой аудит их
коэффициентов обнаруживает общее кольцо
\begin{equation}
 \mathcal A_{24}=\mathbb Z[1/24][\Pi,\Pi^{-1}].
\end{equation}
Все рациональные коэффициенты имеют знаменатели, делящие $24$. Формулы с
$\alpha$ также остаются в поле частных $\operatorname{Frac}(\mathcal A_{24})$,
если подставить объявленное выражение $\alpha=1/S_{vac}(\pi)$.

Единственное исключение среди отобранных коротких формул ---
$\delta_{CKM}=\pi/e$: число Эйлера не порождается этим $\pi$-кольцом.

\section{Универсальный оператор}

\begin{definition}[Оператор масштабирования]
В базисе $|k\rangle$, $k=-4,\ldots,3$, зададим
\begin{equation}
 D|k\rangle=k|k\rangle,
 \qquad
 \Pi=\exp((\log\pi)D).
\end{equation}
Тогда $\Pi|k\rangle=\pi^k|k\rangle$.
\end{definition}

Для диагональных матриц коэффициентов $C_f^{num}$ и $C_f^{den}$ определим
спектральный адрес
\begin{equation}
 \mathcal O_f
 =\frac{\Tr(C_f^{num}\Pi)}{\Tr(C_f^{den}\Pi)}.
\end{equation}
Эта одна конструкция точно воспроизводит все одиннадцать проверенных формул.
Например,
\begin{align}
 \alpha_s&=\frac{1}{6+\pi^2/4},\\
 \frac{m_b}{m_p}&=\pi+\frac43,\\
 \frac{m_d}{m_e}&=\pi^2-1,\\
 \Omega_{dm}&=\pi^{-1}-\frac12\pi^{-2}.
\end{align}

\begin{proposition}[Общая спектральная компрессия]
Одиннадцать сильных $\pi$-формул RPFT допускают представление как отношения
следов одного восьмимерного оператора масштабирования с коэффициентами из
$\mathbb Z[1/24]$.
\end{proposition}

Это утверждение проверено машинно с ошибкой менее $10^{-13}$.

\section{Почему возникает число 24}

Дробные коэффициенты имеют особенно простую интерпретацию как нормированные
ранги в $24$-мерном пространстве:
\begin{equation}
 \frac14=\frac6{24},\qquad
 \frac13=\frac8{24},\qquad
 \frac12=\frac{12}{24},\qquad
 \frac23=\frac{16}{24},\qquad
 \frac34=\frac{18}{24}.
\end{equation}
Число $24$ одновременно является размерностью присоединённого представления
$SU(5)$ и знаменателем Casimir-поправки в $S_{vac}$. Поэтому возникает
проверяемая гипотеза: рациональные части формул могут быть нормированными
следами проекторов в едином $24$-мерном внутреннем секторе, а степени $\pi$ ---
собственными весами оператора масштабирования.

\begin{protocol}[Следующий содержательный тест]
Для каждого наблюдаемого коэффициентная матрица $C_f$ должна быть выведена
только из представления частицы, рангов проекторов, характеров голономии и
нормированного следа. Запрещается выбирать её после просмотра массы или
константы связи.
\end{protocol}

\section{Граница результата}

Общий оператор пока создаёт только степени $\pi$. Для одиннадцати наблюдаемых
всё ещё потребовались $22$ отдельные коэффициентные матрицы и $34$ ненулевых
коэффициентных слота. Следовательно, найдено единое пространство адресов, но
не найден закон, назначающий адрес частице.

\begin{nogocriterion}[Запрет диагональной магии]
Если коэффициентные матрицы нельзя вывести из квантовых чисел до сравнения с
данными, операторное представление является точной перекодировкой атласа, а не
физическим объяснением. Положительный статус возникает только после вывода
минимум двух новых адресов одним и тем же правилом.
\end{nogocriterion}

Численный протокол сохранён в
\texttt{s2t\_pi\_spectral\_address\_operator\_results.json} и воспроизводится
программой \texttt{s2t\_pi\_spectral\_address\_operator\_audit.py}.

\end{document}